\documentclass[11pt, a4paper]{article}

\usepackage[T1]{fontenc}
\usepackage[utf8]{inputenc}
\usepackage{tgpagella}
\usepackage[spanish, es-noshorthands]{babel}
\usepackage{amsmath, amssymb}
\usepackage[table, xcdraw]{xcolor} % Soluciona el error de \rowcolor
\usepackage[most]{tcolorbox}
\usepackage{tabularx}
\usepackage[margin=2.5cm]{geometry}
\usepackage{fancyhdr}

\pagestyle{fancy}
\fancyhf{}
\setlength{\headheight}{16pt}
\lhead{\textbf{UCB Sede Tarija} | Ing. Mecatrónica}
\rhead{IMT-221 | Referencia Rápida}
\renewcommand{\headrulewidth}{0.4pt}
\definecolor{ucbblue}{RGB}{0, 51, 102}

\begin{document}

\begin{tcolorbox}[colback=ucbblue!5, colframe=ucbblue, title=\textbf{Comparación Funcional de Configuraciones del BJT}]
    Esta tabla expone \textbf{tendencias funcionales} de pequeña señal, no garantías numéricas universales. Utiliza esta guía para seleccionar la topología adecuada según las necesidades de impedancia y voltaje del diseño.
\end{tcolorbox}

\vspace{0.5cm}

\renewcommand{\arraystretch}{1.8}
% El formato p{3.8cm} y \raggedright solucionan los Overfull y Underfull \hbox
\begin{tabularx}{\textwidth}{|>{\raggedright\arraybackslash}p{3.8cm}|>{\raggedright\arraybackslash}X|>{\raggedright\arraybackslash}X|>{\raggedright\arraybackslash}X|}
\hline
\rowcolor{ucbblue!20}
\textbf{Propiedad} & \textbf{Emisor Común (CE)} & \textbf{Colector Común (CC) / Seguidor} & \textbf{Base Común (CB)} \\ \hline
\textbf{Terminal de Entrada} & Base & Base & Emisor \\ \hline
\textbf{Terminal de Salida} & Colector & Emisor & Colector \\ \hline
\textbf{Terminal Común (Tierra AC)} & Emisor & Colector & Base \\ \hline
\textbf{Tendencia de Ganancia de Tensión} & Magnitud potencialmente alta ($|A_v| \gg 1$) & Cercana a la unidad ($A_v \approx 1$) & Potencialmente alta ($A_v > 1$) \\ \hline
\textbf{Relación de Fase ($v_o$ vs $v_i$)} & Invertida ($180^\circ$) & No invertida ($0^\circ$) & No invertida ($0^\circ$) \\ \hline
\textbf{Tendencia de Resistencia de Entrada ($R_{in}$)} & Moderada & \textbf{Alta} (minimiza carga en la fuente) & \textbf{Baja} \\ \hline
\textbf{Tendencia de Resistencia de Salida ($R_{out}$)} & Moderada a Alta & \textbf{Baja} (ideal para manejar cargas pesadas) & Alta a Moderadamente Alta \\ \hline
\textbf{Rol Típico en Diseño} & Etapa principal de ganancia de voltaje. & \textbf{Buffer} / Transformador de impedancia. & Interfaz de corriente o diseño de alta frecuencia. \\ \hline
\end{tabularx}

\vspace{1cm}

\subsubsection*{Notas de Carga Multietapa}
\begin{itemize}
    \item La ganancia total de un amplificador de múltiples etapas \textbf{NO} es el producto simple de las ganancias aisladas de cada etapa.
    \item La resistencia de entrada ($R_{in}$) de la segunda etapa actúa como la resistencia de carga ($R_L$) externa para la primera etapa, atenuando su ganancia real en el circuito integrado.
    \item Por ello, el \textbf{Colector Común} es inmensamente valioso a pesar de no amplificar voltaje: su alta $R_{in}$ no afecta la ganancia de la etapa previa, y su baja $R_{out}$ permite alimentar cargas de bajos ohmios sin que la tensión se desplome.
\end{itemize}

\end{document}
